\subsection{Stati coerenti}
Sovrapposizione di stati di oscillatore armonico \( \Rightarrow \) pacchetti d'onda ad indeterminazione minima posizione-impulso che non si allargano perch\'e sono soggetti a potenziale. Definisco questi stati in modo che 
\[
\hat{a} \ket{z} = z \ket{z}, \quad z \in \mathbb{C}, \hat{a} \text{ operatore distruzione},
\]
\[
\mathcal{H} = \frac{\hat{p}^2}{2m} + \frac{1}{2} m\omega^2 \hat{x}^2,
\]
\[
\ket{z} = \sum_n \ket{n}\braket{n|z}.
\]
Ricordo che \( N \ket{n} = n \ket{n}, N = \hat{a}^\dag \hat{a} \), quindi
\begin{align*}
\hat{a} \ket{z} &= \underbrace{\sum_{n = 0}^\infty \hat{a} \ket{n} \braket{n|z}}_{\text{vuoto annichilato}} = \sum_{n = 1}^\infty \sqrt{n} \ket{n - 1} \braket{n|z} \\
&= \sum_{n = 0}^\infty \sqrt{n + 1} \ket{n} \braket{n + 1|z}.
\end{align*}
Allora la condizione di autostato \`e data da 
\[
z \braket{n|z} = \sqrt{n + 1} \braket{n + 1|z}
\]
\[
\braket{n - 1|z} = \sqrt{n} \braket{n|z} \frac{1}{z}
\]
\begin{align*}
\braket{n|z} &= \frac{z}{\sqrt{n}} \braket{n - 1|z} = \\
&= \frac{z^2}{\sqrt{n}\sqrt{n - 1}} \braket{n - 2|z} = \\
&= \frac{z^n}{\sqrt{n!}} \braket{0|z}
\end{align*}
quindi scrivo \( \ket{z} \) come
\begin{align*}
\ket{z} &= \sum_{n = 0}^\infty \ket{n} \frac{z^n}{\sqrt{n!}} \braket{0|z} = \sum_{n = 0}^\infty \frac{(\hat{a}^\dag)^n}{\sqrt{n!}\sqrt{n!}} \ket{0} z^n \braket{0|z} = \\
&= \sum_{n = 0}^\infty \frac{(z\hat{a}^\dag)^n}{n!} \ket{0} \braket{0|z} = e^{z\hat{a}^\dag} \ket{0} \braket{0|z}.
\end{align*}
Impongo che \( \braket{z|z} = 1 \), allora 
\[
1 = |\braket{0|z}|^2 \sum_{n, m} \braket{m|n} \frac{z^n {z^*}^m}{\sqrt{n!m!}} = \sum_{n = 0}^\infty \frac{|z|^{2n}}{n!} |\braket{0|z}|^2 = e^{|z|^2} |\braket{0|z}|^2
\]
\[
\braket{0|z} = e^{-\frac{1}{2}|z|^2}
\]
\[
\boxed{\ket{z} = e^{-\frac{1}{2}|z|^2} e^{z\hat{a}^\dag} \ket{0}},
\]
questa \`e la forma degli \textbf{stati coerenti}.

I valori medi sono 
\[
\braket{z|\hat{x}|z} = \sqrt{\frac{\hbar}{2m\omega}} \braket{z|(\hat{a} + \hat{a}^\dag)|z} = \sqrt{\frac{\hbar}{2m\omega}} (z + z^*) = \sqrt{\frac{2\hbar}{m\omega}} \, \text{Re} \,  z
\]
\[
\braket{z|\hat{p}|z} = -i\sqrt{\frac{m\hbar\omega}{2}}\braket{z|\hat{a} - \hat{a}^\dag|z} = \sqrt{2m\omega\hbar} \, \text{Im} \, z
\]
e quindi sono valori oscillanti. Inoltre
\[
\braket{\hat{x}} = 2 \, \text{Re}z = \Delta \hat{x},
\]
\[
\braket{\hat{p}} = 2 \, \text{Im}z = \Delta\hat{p}.
\]

Gli scarti quadratici
\begin{align*}
\braket{z|\Delta^2 \hat{x}|z} &= \braket{z|(\hat{x} - \braket{\hat{x}})^2|z} = \\
&= \frac{\hbar}{2m\omega} \braket{z|((\hat{a} + \hat{a}^\dag) - (\braket{\hat{a}} + \braket{\hat{a}^\dag}))((\hat{a} + \hat{a}^\dag) - (\braket{\hat{a}} + \braket{\hat{a}^\dag}))|z} = \\
&= \frac{\hbar}{2m\omega} \braket{z|(\hat{a} + z^* - z - z^*)(z + \hat{a}^\dag - z - z^*)|z} = \\
&= \frac{\hbar}{2m\omega} \braket{z|(\hat{a} - z)(\hat{a}^\dag - z^*)|z} = \\
&= 0 + \frac{\hbar}{2m\omega} \braket{z|[\hat{a} - z, \hat{a}^\dag - z^*]|z} = \\
&= \frac{\hbar}{2m\omega}, \\
\braket{z|\Delta^2 \hat{p}|z} &= \frac{\hbar m\omega}{2}
\end{align*}
danno l'indeterminazione minima 
\[
\boxed{\Delta^2 \hat{x} \Delta^2 \hat{p} = \frac{\hbar^2}{4}}.
\]
Ai tempi successivi, poich\'e \( \hat{a}(t) = e^{-i\omega t}\hat{a} \), l'indeterminazione si conserva, mentre i valori medi oscillano.

Per due stati coerenti 
\[
\braket{z_1|z_2} = \sum_{n = 0}^\infty \frac{(z_1^* z_2)^n}{n!} e^{-\frac{1}{2}|z_1|^2} e^{\frac{1}{2}|z_2|^2} = e^{z_1^* z_2} e^{-\frac{|z_1|^2 - |z_2|^2}{2}},
\]
ad esempio, per \( \braket{z|-z} = e^{-2|z|^2} \), cio\`e ho un decadimento esponenziale di \( \braket{z|-z} \), ma poich\'e 
\[
|z|^2 = (\text{Re}z)^2 + (\text{Im}z)^2 = \frac{\braket{\hat{p}}^2}{4(\Delta\hat{p})^2} + \frac{\braket{\hat{x}}^2}{4(\Delta\hat{x})^2},
\]
\[
\braket{z|-z} = e^{2 \left( \frac{1}{4} + \frac{1}{4} \right)} = e.
\]